\documentclass[11pt,a4paper,spanish]{book}
\usepackage{estilo_unir-1}
\usepackage{apacite}

\usepackage{hyperref}
\numberwithin{equation}{chapter}
\numberwithin{figure}{chapter}

%---------------------------
%título del trabajo y autor
%---------------------------
\title{ProcureWatch Analytics: Planteamiento Metodológico}
\titulacion{Máster Universitario en Análisis y Visualización de Datos Masivos}
\author{Sebastián Alfonso Correa y Saturía María Hurtado Álvarez}
\date{22 de abril de 2026}
\director{Marlon Cárdenas Bonett}
\nombreciudad{Madrid}

%---------------------------
%---------------------------

\begin{document}
\renewcommand{\listfigurename}{Índice de Ilustraciones}
\renewcommand{\listtablename}{Índice de Tablas}
\renewcommand{\contentsname}{Índice de Contenidos}
\renewcommand{\figurename}{Figura}
\renewcommand{\tablename}{Tabla}

\maketitle

\frontmatter
\tableofcontents
\listoffigures
\listoftables

\chapter{Prefacio}

Este documento contiene el planteamiento metodológico detallado del proyecto ProcureWatch Analytics. Describe las fases de trabajo, organización del equipo, cronograma, entregables y criterios de calidad para la ejecución del Trabajo Fin de Máster.

\mainmatter

%------------------------------------------
%------------------------------------------
\chapter{Planteamiento Metodológico}

\section{Enfoque General}

El proyecto sigue una \textbf{metodología iterativa adaptada de CRISP-DM} para la componente analítica, combinada con \textbf{desarrollo incremental por módulos} para el prototipo. Esta aproximación permite validación temprana de hipótesis, integración progresiva de módulos y feedback continuo.

\section{Fases de Trabajo}

\subsection{Fase 0: Configuración del Entorno (Semanas 1-2)}

\textbf{Objetivo:} Preparar infraestructura técnica y recursos formativos.

\textbf{Actividades:}
\begin{itemize}
\item Instalación local: Python, Docker, Ollama, PostgreSQL, Neo4j Community
\item Setup de repositorio GitHub con estructura de carpetas
\item Cursos introductorios: Python para DA, Fundamentos de IA
\end{itemize}

\textbf{Entregables:} Entorno funcional, repositorio inicializado, documentación de setup.

\textbf{Responsables:} Ambos
\textbf{Duración:} 2 semanas

---

\subsection{Fase 1: Comprensión del Dominio (Semanas 3-4)}

\textbf{Objetivo:} Entender estructura del ciclo contractual español y caracterizar calidad de datos.

\textbf{Actividades:}
\begin{itemize}
\item Exploración estadística del dataset BOE 2014-2024
\item Análisis comparativo PLACE vs BOE
\item Documentación del ciclo contractual español
\item Revisión de red flags en literatura (OCP 2024, DIGIWHIST)
\item Elaboración de matriz: Variables OCDS $\times$ Indicadores de red flags
\end{itemize}

\textbf{Entregables:}
\begin{enumerate}
\item Informe de análisis exploratorio (EDA report) con gráficos y estadísticas descriptivas
\item Documento técnico de mapeo: campos OCDS $\rightarrow$ indicadores de red flags
\item Lista priorizada de red flags a implementar (mínimo 15)
\item Presentación interna con hallazgos y recomendaciones
\end{enumerate}

\textbf{Responsables:} Ambos (complementarios)
\textbf{Duración:} 2 semanas

---

\subsection{Fase 2: Preparación de Datos (Semanas 5-6)}

\textbf{Objetivo:} Construir pipeline de datos limpio, normalizado y disponible para análisis.

\textbf{Actividades:}
\begin{itemize}
\item Implementación del pipeline ETL: descarga, validación, limpieza, normalización OCDS
\item Resolución de inconsistencias: duplicados, fechas incoherentes, importes negativos, NIFs malformados
\item Enriquecimiento: integración con datos de registro mercantil (cuando disponible)
\item Carga en PostgreSQL con validaciones post-insert
\item Construcción del grafo empresa-organismo en Neo4j usando Python driver
\end{itemize}

\textbf{Entregables:}
\begin{enumerate}
\item Script pipeline reproducible y documentado (pipeline.py)
\item Base de datos PostgreSQL con esquema OCDS completo: mínimo 97.000 contratos, 2.000+ organismos, 15.000+ empresas
\item Grafo Neo4j operacional con estructura verificable
\item Reporte de calidad post-carga (completitud, coherencia temporal, duplicados)
\end{enumerate}

\textbf{Responsable:} Sebastián (lead pipeline), Saturía (apoyo)
\textbf{Duración:} 2 semanas

---

\subsection{Fase 3: Motor de Red Flags (Semanas 7-9)}

\textbf{Objetivo:} Implementar y validar indicadores de red flags y modelos de scoring.

\textbf{Actividades:}
\begin{itemize}
\item Codificación de 15+ indicadores (RF-01 a RF-06 y variantes) en SQL/Python
\item Implementación de scoring compuesto: base score + bonus/penalty por métricas de red
\item Comparativa de tres enfoques: reglas OCDS explícitas vs. Isolation Forest vs. Positive-Unlabeled Learning
\item Validación parcial: análisis de precisión sobre casos documentados (contratos menores con licitación única verificables)
\item Ajuste de umbrales y pesos según resultados
\end{itemize}

\textbf{Entregables:}
\begin{enumerate}
\item Módulo de red flags implementado (red\_flags.py) con documentación de cada indicador
\item Scoring por contrato y por adjudicatario disponible en BD
\item Comparativa técnica: tablas y gráficos de rendimiento de tres enfoques
\item Ranking de red flags por frecuencia de activación y correlación con casos documentados
\end{enumerate}

\textbf{Responsable:} Sebastián (lead)
\textbf{Duración:} 3 semanas

---

\subsection{Fase 4: Análisis de Redes y NLP (Semanas 10-12)}

\textbf{Objetivo:} Análisis relacional (comunidades, centralidades) y extracción documental.

\textbf{Subactividad 4a — Análisis de Redes (Saturía):}

\begin{itemize}
\item Ejecución de algoritmo Leiden en grafo empresa-organismo: detección de comunidades
\item Cálculo de métricas de centralidad: degree, betweenness, eigenvector
\item Detección de patrones anómalos: adjudicatarios con centralidad extrema, comunidades mal conectadas
\item Correlación entre métricas de red y scores de red flags
\end{itemize}

\textbf{Subactividad 4b — NLP y OCR (Saturía):}

\begin{itemize}
\item OCR de pliegos y resoluciones con Docling
\item Extracción de entidades con spaCy (organismos, importes, plazos, proveedores)
\item Construcción de corpus indexado en Qdrant (vector database)
\item Implementación de motor RAG híbrido (dense retrieval + sparse retrieval)
\end{itemize}

\textbf{Entregables:}
\begin{enumerate}
\item Informe de análisis de comunidades: número, tamaño, modularidad, patrones atípicos
\item Corpus documental procesado y indexado (~5.000 documentos)
\item Motor RAG funcionando con métricas RAGAS documentadas
\item Visualización: grafo interactivo con comunidades coloreadas, centralidades etiquetadas
\end{enumerate}

\textbf{Responsable:} Saturía (lead ambas subactividades)
\textbf{Duración:} 3 semanas

---

\subsection{Fase 5: Sistema Multiagente y Front-end (Semanas 13-14)}

\textbf{Objetivo:} Integrar todos los módulos en plataforma cohesiva.

\textbf{Actividades:}
\begin{itemize}
\item Integración de agentes con LangGraph: orquestación de pipeline $\rightarrow$ red flags $\rightarrow$ redes $\rightarrow$ NLP
\item Implementación de front-end con Streamlit + Plotly + Sigma.js
\item Componentes: grafo interactivo, panel de indicadores, timeline, ficha explicativa
\item Conexión backend (APIs FastAPI) $\leftrightarrow$ front-end
\item Tests end-to-end: flujos completos de usuario
\end{itemize}

\textbf{Entregables:}
\begin{enumerate}
\item Plataforma integrada y funcional (demo ejecutable)
\item Interfaz web con componentes interactivos
\item API REST documentado (OpenAPI/Swagger)
\item Guía de usuario técnico
\end{enumerate}

\textbf{Responsables:} Ambos (Sebastián back-end, Saturía front-end)
\textbf{Duración:} 2 semanas

---

\subsection{Fase 6: Evaluación y Refinamiento (Semanas 15-16)}

\textbf{Objetivo:} Validación rigurosa del prototipo con métricas cuantitativas y cualitativas.

\textbf{Actividades:}
\begin{itemize}
\item Ejecución de suite de pruebas: pipeline quality, red flags precision/recall, RAG metrics, usability SUS
\item Análisis detallado de 10-15 casos de riesgo: comparación scoring sistema vs. realidad documentada
\item Refinamiento de pesos, umbrales, modelos basado en resultados
\item Documentación de limitaciones, gaps, roadmap futuro
\end{itemize}

\textbf{Entregables:}
\begin{enumerate}
\item Informe de evaluación completo con todas las métricas
\item Análisis de casos: fichas técnicas de 10-15 contratos analizados
\item Documento de lecciones aprendidas y mejoras futuras
\end{enumerate}

\textbf{Responsables:} Ambos
\textbf{Duración:} 2 semanas

---

\subsection{Fase 7: Documentación y Cierre (Semanas 17-18)}

\textbf{Objetivo:} Redacción de memoria TFM, documentación técnica y preparación de defensa.

\textbf{Actividades:}
\begin{itemize}
\item Redacción de memoria TFM siguiendo estructura UNIR: introducción, estado del arte, objetivos, metodología, resultados, conclusiones
\item Documentación técnica exhaustiva (architecture, APIs, ejemplos)
\item Preparación de repositorio GitHub público con código, datos (subset) y documentación
\item Ensayos de presentación y defensa
\end{itemize}

\textbf{Entregables:}
\begin{enumerate}
\item Memoria TFM completa y editada (~100-150 páginas)
\item Repositorio GitHub con código fuente, documentación, ejemplos reproducibles
\item Presentación para defensa (~20-25 minutos)
\end{enumerate}

\textbf{Responsables:} Ambos
\textbf{Duración:} 2 semanas

---

\section{Organización del Equipo}

El equipo consta de \textbf{dos investigadores} con roles complementarios:

\subsection{Sebastián Alfonso Correa}

\textbf{Especialidad:} Pipeline de datos, modelado analítico, machine learning, scoring

\textbf{Responsabilidades principales:}
\begin{itemize}
\item Fase 2: Diseño e implementación del pipeline ETL (lead)
\item Fase 3: Desarrollo del motor de red flags (lead)
\item Fase 5: Back-end, APIs, orquestación de agentes
\item Apoyo transversal en validación de datos
\end{itemize}

\textbf{Entregables específicos:}
\begin{itemize}
\item Pipeline reproducible (clean\_data.py, load\_database.py)
\item Módulo de red flags (red\_flags.py, scoring.py)
\item API backend (FastAPI, main.py)
\item Documentación técnica de arquitectura
\end{itemize}

---

\subsection{Saturía María Hurtado Álvarez}

\textbf{Especialidad:} Análisis de redes, NLP, procesamiento de documentos, visualización

\textbf{Responsabilidades principales:}
\begin{itemize}
\item Fase 1: Análisis exploratorio y matriz de mapeo (co-lead)
\item Fase 4: Análisis de grafos y redes (lead)
\item Fase 4: NLP, OCR y motor RAG (lead)
\item Fase 5: Front-end y componentes de visualización
\item Apoyo en evaluación cualitativa de casos
\end{itemize}

\textbf{Entregables específicos:}
\begin{itemize}
\item Análisis exploratorio (eda\_report.py)
\item Módulo de análisis de redes (network\_analysis.py)
\item Pipeline NLP (nlp\_pipeline.py, rag\_engine.py)
\item Front-end interactivo (streamlit\_app.py)
\item Visualizaciones (grafo interactivo, dashboards)
\end{itemize}

---

\subsection{Sincronización y Coordinación}

\begin{itemize}
\item \textbf{Reuniones semanales de equipo:} revisión de progreso, blockers, decisiones de diseño
\item \textbf{Hito de revisión de Fase 2/3:} validación de datos y scoring antes de avanzar a redes
\item \textbf{Integración progresiva:} módulos se desarrollan en paralelo pero con interfases claras (BD, APIs)
\item \textbf{Testing compartido:} cada módulo debe pasar tests del otro antes de mergear a rama principal
\end{itemize}

---

\section{Entregables por Fase}

\begin{table}[h]
\centering
\caption{Resumen de fases, entregables y responsables}
\label{tab:fases_entregables}
\begin{tabular}{|p{0.8cm}|p{1.5cm}|p{3cm}|p{1.5cm}|p{1cm}|}
\hline
\textbf{Fase} & \textbf{Nombre} & \textbf{Entregables clave} & \textbf{Responsable} & \textbf{Sem} \\
\hline
0 & Configuración & Entorno, repo, cursos & Ambos & 2 \\
\hline
1 & Análisis explorat. & EDA report, matriz vars-flags & Ambos & 2 \\
\hline
2 & Preparación datos & Pipeline, BD, grafo & Sebastián & 2 \\
\hline
3 & Motor red flags & Módulo RF, scoring, comparativa & Sebastián & 3 \\
\hline
4a & Análisis redes & Comunidades, patrones & Saturía & 3 \\
\hline
4b & NLP/OCR & Corpus, motor RAG & Saturía & 3 \\
\hline
5 & Front-end integrado & Plataforma, APIs, interfaz & Ambos & 2 \\
\hline
6 & Evaluación & Informe eval., casos, roadmap & Ambos & 2 \\
\hline
7 & Cierre & Memoria, repo, presentación & Ambos & 2 \\
\hline
\end{tabular}
\end{table}

\textbf{Duración total:} 18 semanas (~4.5 meses)

---

\section{Criterios de Calidad y Validación}

\subsection{Criterios técnicos por módulo}

\begin{itemize}
\item \textbf{Pipeline:} completitud OCDS $>$ 90\%, coherencia temporal $>$ 98\%
\item \textbf{Red Flags:} precisión $>$ 70\%, estabilidad $\pm$ 10\%
\item \textbf{Redes:} modularidad $Q > 0.30$, correlación con flags $\geq 0.45$
\item \textbf{NLP/RAG:} Faithfulness $>$ 0.80, Answer Relevancy $>$ 0.75, Context Recall $>$ 0.70
\item \textbf{Front-end:} SUS $>$ 68, tiempo medio caso $<$ 5 minutos
\end{itemize}

\subsection{Criterios de reproducibilidad}

\begin{itemize}
\item Código completamente documentado y versionado en GitHub
\item Scripts ejecutables con ejemplos reproducibles
\item Dataset base (subset de BOE) incluido en repo o con instrucciones de descarga
\item Docker compose para despliegue local
\end{itemize}

\subsection{Criterios de pertinencia académica}

\begin{itemize}
\item \textbf{Estado del arte:} mínimo 10 referencias académicas en formato APA por bloque temático
\item \textbf{Contribución clara:} gap identificado y cerrado vs. literatura existente
\item \textbf{Rigor metodológico:} métricas estándar de la literatura adaptadas a contexto español
\end{itemize}

---

\chapter{Gestión de Riesgos y Contingencias}

\section{Riesgos Identificados}

\subsection{Riesgo 1: Calidad de datos inferior a expectativas}

\textbf{Probabilidad:} Media \textbf{|} \textbf{Impacto:} Alto

\textbf{Mitigación:}
\begin{itemize}
\item Fase 1: análisis exhaustivo de completitud antes de comprometer Phase 2
\item Identificar fuentes complementarias (registros mercantiles, ROLECE)
\item Documentar limitaciones de calidad en reporte final
\end{itemize}

\subsection{Riesgo 2: Alcance técnico demasiado ambicioso}

\textbf{Probabilidad:} Media \textbf{|} \textbf{Impacto:} Alto

\textbf{Mitigación:}
\begin{itemize}
\item Priorizar MVP (red flags básicos + grafo) sobre funcionalidades avanzadas
\item Hito de revisión en Fase 3: evaluar scope vs. tiempo disponible
\item Documentar funcionalidades descartadas como trabajo futuro
\end{itemize}

\subsection{Riesgo 3: Dificultades en integración multiagente}

\textbf{Probabilidad:} Media \textbf{|} \textbf{Impacto:} Medio

\textbf{Mitigación:}
\begin{itemize}
\item Interfaces claras entre módulos (BD + APIs documentadas)
\item Testing early and often
\item Slack de 1 semana en Fase 5 para integración
\end{itemize}

---

\appendix
\chapter{Estructura de Carpetas del Proyecto}

\begin{verbatim}
procurewatch-analytics/
├── data/
│   ├── raw/          # Datos descargados sin procesar
│   ├── processed/    # Datos limpios y normalizados
│   └── external/     # Datos complementarios
├── src/
│   ├── pipeline/     # Pipeline ETL
│   ├── red_flags/    # Motor de red flags
│   ├── networks/     # Análisis de redes
│   ├── nlp/          # NLP y RAG
│   └── frontend/     # Interfaz web
├── notebooks/       # Análisis exploratorio
├── tests/          # Suite de pruebas
├── docs/           # Documentación
├── docker-compose.yml
├── requirements.txt
└── README.md
\end{verbatim}

\chapter{Reuniones y Hitos Clave}

\section{Calendario de Reuniones}

\begin{itemize}
\item \textbf{Semanal (lunes 10:00):} Reunión de sincronización equipo (30 min)
\item \textbf{Bi-semanal (viernes 16:00):} Review con director de tesis (1 hora)
\item \textbf{Fin de fase:} Presentación de hitos y entregables (1.5 horas)
\end{itemize}

\section{Hitos Críticos}

\begin{enumerate}
\item \textbf{Semana 2:} Entorno completamente funcional
\item \textbf{Semana 4:} BD poblada con datos limpios y validados
\item \textbf{Semana 9:} Motor de red flags funcional (MVP)
\item \textbf{Semana 12:} Análisis de redes y NLP completados
\item \textbf{Semana 14:} Plataforma integrada y funcional
\item \textbf{Semana 16:} Evaluación completa y documentación
\item \textbf{Semana 18:} Memoria TFM y presentación listos para defensa
\end{enumerate}

\end{document}
